\documentclass[twocolumn,10pt]{article}

\usepackage[utf8]{inputenc}
\usepackage{amsmath}
\usepackage{amssymb}
\usepackage{booktabs}
\usepackage{microtype}
\usepackage{geometry}
\geometry{a4paper, margin=0.75in}

\title{\textbf{An Analytical Framework for Multi-Variable Non-Linear Damage Superposition in Bounded Discrete Tactical Environments}}
\author{\textbf{BotAn} \\ \small Brown Dust II Tactical Compendium}
\date{\today}

\begin{document}
\raggedbottom

\maketitle

\begin{abstract}
Predicting discrete state-vector outputs in complex tactical computational environments remains an analytical challenge when underlying source code is obfuscated. This paper presents a phenomenological framework that maps the conditional mechanics of the Brown Dust II engine into a unified, continuous tensor equation. Rather than relying on fragmented imperative control-flow logic, a declarative model is constructed from bottom-up empirical constraints. By representing elemental attributes as orthogonal property vectors and mapping modifiers via multidimensional linear transformations, a mathematically rigorous alternative to traditional algorithmic scripting is provided. Empirical validation across live server environments demonstrates absolute convergence ($R^2 = 1.0$), proving that observational modeling can perfectly reconstruct hidden system rules.
\end{abstract}

\section{Introduction}
In discrete tactical simulation environments, the determination of state transitions—specifically, the calculation of kinetic and ethereal energy transfer between opposing nodes—is a critical computational process. Modern turn-based architectures have evolved past elementary linear scaling systems, introducing multi-layered, interdependent modifier pools. While these complex mechanics enrich the strategic depth of the environment, they frequently introduce severe optimization challenges for system analysis.

Within the specific framework of the Brown Dust II tactical engine, automated data routing and character configuration require a precise understanding of how discrete attributes scale under active status modifiers. In practical execution, players and developers alike encounter a high-dimensional combinatorial optimization problem when configuring team sequencing matrices. Because the underlying source code of the execution engine remains proprietary and obfuscated, the precise topological structure of these interactions cannot be verified through direct code inspection.

This systemic lack of transparency presents a distinct analytical hurdle. When multiple status-altering modifiers—such as offensive amplifications (e.g., Attack or Magic Attack pooling), temporal damage vectors (Chains), and structural vulnerabilities—interact simultaneously, predicting the exact output state becomes non-trivial. Compounding this challenge is the presence of hard operational boundary limits, such as saturation caps applied to entity vitality metrics, which inject non-linear constraints into otherwise predictable algebraic equations.

To resolve these ambiguities, this study adopts a black-box phenomenological approach to reverse-engineer and formalize the global damage system. By conducting systematic, isolated live-server experiments, individual systemic rules are identified, evaluated, and translated from sequential control-flow constraints into a unified mathematical syntax. It is demonstrated that apparent algorithmic discrepancies—such as the transition between additive pooling and multiplicative compounding—can be elegantly modeled using linear algebra and orthogonal basis vector operations.

The primary contributions of this paper are threefold:
\begin{enumerate}
  \item A modular, bottom-up mathematical derivation of individual modifier behaviors constructed entirely from empirical observation.
  \item The formulation of a unified master tensor equation that bypasses the need for imperative if/else control logic.
  \item Computational verification proving that the proposed declarative model achieves absolute predictive convergence with live system executions.
\end{enumerate}

The structural composition of the remainder of this document is organized as follows. Section~\ref{sec:derivation} details the phase-by-phase axiomatic construction of individual scaling components and establishes the mathematical rationale behind additive buff stacking constraints. Section~\ref{sec:master} synthesizes these isolated components into a single, comprehensive master equation. Section~\ref{sec:empirical} presents empirical validation scenarios derived from live server readouts to verify the predictive robustness of the model. Finally, Section~\ref{sec:optimization} leverages the formalized framework to outline the absolute mathematical boundaries of team matrix optimization.

% =========================================================================
% SECTION II: SYSTEM MODELING AND CORE VECTOR MECHANICS
% =========================================================================
\section{System Modeling and Core Vector Mechanics}
\label{sec:derivation}

To construct a verifiable mathematical model of the calculation pipeline, the game's underlying execution engine is analyzed by gradually introducing systemic complexities. This progressive approach ensures that the mathematical framework remains tightly aligned with observable live-server behavior.

\subsection{The Elementary Scalar Model}
The most intuitive baseline model assumes a standard physical interaction where a source unit executes an action against a single target tile occupied by a neutral target unit. In this idealized state, devoid of active status modifiers, elemental properties, or defensive mitigation mechanics, the primary baseline damage output ($D_0$) can be expressed as a basic scalar product:
\begin{equation}
D_0 = \text{ATK} \times M
\label{eq:scalar_base}
\end{equation}
where $\text{ATK}$ represents the flat physical attack attribute of the source unit, and $M$ denotes the scaling multiplier dictated by the active skill. While Equation~\ref{eq:scalar_base} successfully accounts for elementary, single-variable physical interactions, it quickly fails when confronted with more sophisticated unit skill sets.

\subsection{The Attribute Dependency Conflict}
Empirical testing across diverse unit skill sets reveals that the calculation engine does not utilize a single, uniform attribute for output evaluations. Instead, independent character and costume skill designs shift their scaling dependencies dynamically across entirely distinct statistical pools:
\begin{enumerate}
    \item \textbf{Magical Payloads:} Certain skill sets completely swap physical parameters for a magical alternative, forcing the baseline calculation to rely on $\text{MATK} \times M$
    \item \textbf{Vitality Scaling:} Advanced tank-type units leverage alternative skill designs that bypass offensive statistics entirely, generating an output value scaled by the source unit's own maximum health points ($\text{HP} \times M$).
    \item \textbf{Simultaneous Multi-Attribute Scaling:} The most critical analytical conflict occurs with hybrid skill sets that track multiple independent metrics simultaneously. For example, specific protective mechanics compute an execution value based on the concurrent summation of both the source unit's current Health Points and their active protective barrier layers (Energy Guard, $\text{EG}$) on a localized target tile.
\end{enumerate}
Attempting to model these fluid, multi-variable dependencies using standard scalar mathematics would require a disjointed set of nested conditional control statements. To maintain a clean, declarative mathematical framework that handles single-attribute and multi-attribute calculations uniformly, a structural shift from scalar spaces to vector spaces is required.

\subsection{Vector Space Formalization and Operational Constraints}
To elegantly resolve the multi-attribute dependencies identified above, alongside interactive mechanics that scale dynamically based on target unit metrics (such as target-punishment or executioner scaling properties), the interaction state is compiled into a unified, seven-dimensional column vector, designated as the Source-Target Stat Vector, $\vec{v} \in \mathbb{R}^7$:
\begin{equation}
\vec{v} = \begin{pmatrix} 
\text{ATK}_{\text{src}} \\ 
\text{MATK}_{\text{src}} \\ 
\text{HP}_{\text{src}} \\ 
\text{EG}_{\text{src}} \\ 
\text{ATK}_{\text{tgt}} \\ 
\text{MATK}_{\text{tgt}} \\ 
\text{HP}_{\text{tgt}} 
\end{pmatrix}
\end{equation}
This vertical coordinate space maps the primary offensive parameters alongside vitality pools for both the source entity and the target entity simultaneously.

At the database level, the system enforces strict structural boundaries on the elements of $\vec{v}$ based on the character templates assigned to the respective units. Because any individual unit asset is hardcoded to belong exclusively to either a physical or a magical archetype, the active offensive parameters within the vector space must satisfy a pair of decoupled zero-product boundary conditions:
\begin{equation}
v_1 \cdot v_2 = 0 \quad (\text{ATK}_{\text{src}} \cdot \text{MATK}_{\text{src}} = 0)
\label{eq:exclusivity_src}
\end{equation}
\begin{equation}
v_5 \cdot v_6 = 0 \quad (\text{ATK}_{\text{tgt}} \cdot \text{MATK}_{\text{tgt}} = 0)
\label{eq:exclusivity_tgt}
\end{equation}
These conditions ensure that if a unit configuration yields active physical parameters ($\text{ATK} \neq 0$), its corresponding magical component is fundamentally driven to zero ($\text{MATK} = 0$), and vice versa. Vitality metrics face no such structural constraints and exist simultaneously across all active unit entities.

To interface with this combined attribute snapshot, the executing skillset maps its specific scaling coefficients into a matching seven-dimensional column vector, designated as the Skill Scaling Vector, $\vec{m} \in \mathbb{R}^7$
\begin{equation}
\vec{m} = \begin{pmatrix} m_1 & m_2 & m_3 & m_4 & m_5 & m_6 & m_7 \end{pmatrix}^{\top}
\end{equation}
All individual scaling coefficients within the Skill Scaling Vector are bound by a strict non-negativity constraint, such that $m_i \ge 0$ for all element index positions $i \in \{1, \dots, 7\}$. This structural floor ensures that a baseline skillset execution can only generate constructive or neutral payload components, mathematically isolating all reductive transformations to the external debuff pools and adversarial mitigation layers evaluated later in the pipeline.

To calculate the active interaction payload components across all seven statistical channels simultaneously, the engine executes an element-wise Hadamard product ($\circ$) between the two column vectors, creating an intermediate column Execution Vector, $\vec{e} \in \mathbb{R}^7$:
\begin{equation}
\vec{e} = \vec{v} \circ \vec{m}
\label{eq:hadamard_exec}
\end{equation}
where each individual element is evaluated via scalar multiplication as $e_i = v_i \cdot m_i$. Utilizing the element-wise mapping of the Hadamard product is mathematically mandatory for this system configuration; it ensures that independent statistical properties remain completely isolated before final summation. This architectural choice allows subsequent environmental modifier fields to be mapped onto the system as independent vector operations.

The final, unmitigated base damage output ($D_0$) for the targeted tile is then extracted by evaluating the sum of the resulting execution elements:
\begin{equation}
D_0 = \sum_{i=1}^{7} e_i = \sum_{i=1}^{7} v_i \cdot m_i
\label{eq:hadamard_sum}
\end{equation}

This formulation successfully unifies the underlying system mechanics under a single declarative rule. For a standard physical interaction, Equations~\ref{eq:exclusivity_src} and \ref{eq:exclusivity_tgt} force the magical elements of the active unit profiles to zero. Consequently, the intermediate Execution Vector $\vec{e}$ evaluates to an axis-aligned state containing only the active physical payload components, rendering the final summation in Equation~\ref{eq:hadamard_sum} mathematically equivalent to an elementary scalar model. For hybrid or target-interactive skill sets, the Hadamard framework natively preserves the distinct parallel channels for vitality metrics and target scaling values, mapping out their simultaneous contributions before final aggregation without requiring fragmented conditional code branches.

\subsection{Base Attribute Construction and Level Dependencies}
The individual elements compiled within the Source-Target Stat Vector $\vec{v}$ are not primitive, static constants pulled directly from character sheets. Instead, a structural division exists between persistent, macro-progression attributes controlled by the source entity, transient combat barrier metrics, and external values mirrored from the target state node. 

For the native, persistent source attribute channels $v_i$ where $i \in \{1, 2, 3\}$ (corresponding explicitly to $\text{ATK}_{\text{src}}$, $\text{MATK}_{\text{src}}$, and $\text{HP}_{\text{src}}$ respectively), each element represents the dynamic terminal output of a multi-layered Attribute Aggregation Phase. Live-server empirical research reveals that the engine evaluates a strict hierarchical sequence of level-scaling, flat additions, and global percentage modifiers to construct these final pre-combat stat parameters before grid deployment.

To formalize this persistent pipeline, let $S_{\text{base}}$ represent the initial, native baseline value of a given source attribute channel $i$. The level-dependent growth curve is driven by a linear scaling coefficient, $\alpha_i$, determined explicitly by the category of the target attribute:
\begin{equation}
\alpha_i = 
\begin{cases} 
0.1 & \text{if } i = 1 \text{ or } 2 \quad (\text{ATK}_{\text{src}} \text{ or } \text{MATK}_{\text{src}}) \\ 
0.2 & \text{if } i = 3 \quad (\text{HP}_{\text{src}}) 
\end{cases}
\label{eq:growth_coeff}
\end{equation}
Crucially, while Equation~\ref{eq:growth_coeff} evaluates a non-zero linear progression rate for both offensive channels, the structural boundary constraint established in Equation~\ref{eq:exclusivity_src} dictates that the vacant archetype channel initiates with a baseline property value of $S_{\text{base}} = 0$, rendering its concurrent level aggregation trivially zero.

By utilizing this coefficient alongside the unit's current level ($L$), the absolute flat component accumulation value ($v_{i,\text{flat}}$) aggregates native growth concurrently with flat numerical additions from gear configurations, active or permanent bond systems, and awakening milestones:
\begin{equation}
\begin{split}
v_{i,\text{flat}} = {} & \left( S_{\text{base}} + \alpha_i \cdot L \cdot S_{\text{base}} \right) \\
& + A_{\text{gear\_flat}} + A_{\text{bond\_flat}} + A_{\text{awake\_flat}}
\end{split}
\label{eq:flat_stat_construction}
\end{equation}

Once $v_{i,\text{flat}}$ is established, the engine evaluates a unified multiplicative layer containing all passive, permanent, and equipment-based percentage scaling pools to yield the compound coefficient channel, $v_{i,\text{field}}$:
\begin{equation}
\begin{split}
v_{i,\text{field}} = {} & 1 + \omega_{\text{collection}} + \omega_{\text{gear}} \\
& + \omega_{\text{awakening}} + \omega_{\text{bond}}
\end{split}
\label{eq:field_stat_construction}
\end{equation}
where each $\omega$ term represents the fractional value of its respective percentage bonus pool. The terminal active value loaded into these persistent vector coordinates is subsequently computed as the direct product of these two distinct structural phases:
\begin{equation}
v_i = v_{i,\text{flat}} \times v_{i,\text{field}} \quad \text{for } i \in \{1, 2, 3\}
\label{eq:final_vector_element}
\end{equation}

In contrast to the persistent attribute pipeline, the protective barrier layer channel $v_4$ ($\text{EG}_{\text{src}}$) represents a highly volatile, dynamic instance state rather than an internal progression attribute. Because defensive shield parameters are generated at runtime through diverse external skill variables and non-native scaling references, $v_4$ fundamentally bypasses Equations~\ref{eq:growth_coeff} through \ref{eq:final_vector_element}. Instead, the engine isolates this slot from the macro-progression accumulation layer, populating $v_4$ directly via a real-time snapshot of the absolute active barrier shielding remaining on the source node's grid matrix location at the exact execution frame.

This synthesis successfully concludes the initialization and consolidation of the active attribute parameters. By fully encapsulating all character growth, equipment additions, and passive account progressions directly within the core source coordinates $v_1$ through $v_4$, the system establishes a stable, fully aggregated dataset ready to be processed through the localized skillset execution layers, while elements $v_5$ through $v_7$ are concurrently mirrored from the corresponding target unit's independent aggregation layer at the moment of runtime execution.

\end{document}